\section{Diagnostic Methodologies}
\label{sec:methods}

\subsection{Unified Preprocessing Pipeline}
To ensure rigorous benchmarking, all models consume identical preprocessed inputs derived from a single canonical function. A raw trace $v(0, t)$ is normalized against the incident plateau step $v_{\mathrm{inc}}$:
\begin{equation}
\tilde{v}(t) = \frac{v(0, t)}{v_{\mathrm{inc}}}
\label{eq:v_norm}
\end{equation}
Three parallel feature channels are constructed:
\begin{enumerate}
    \item \textbf{Normalized Voltage $\tilde{v}(t)$:} Encodes baseline plateaus and reflection step levels.
    \item \textbf{Numerical Derivative $d\tilde{v}/dt$:} Implemented via a Savitzky-Golay finite impulse response (FIR) filter (window length 11, polynomial degree 2, first derivative) wrapped in a 1D convolution layer. This sharpens inflection points and isolates high-frequency wave arrival edges while suppressing high-frequency AWGN.
    \item \textbf{Dimensionless Coordinate Channel $t/T_{\mathrm{nom}}$:} Linearly ramps from $0.0$ to $2.5 \cdot (1023/1024)$, providing explicit temporal and spatial reference coordinates.
\end{enumerate}
The resulting input tensor has shape $[3, 1024]$ across batch samples.

\subsection{Conventional TDR Solver (Rule-Based Baseline)}
The conventional baseline implements textbook TDR reasoning with no learned parameters other than a handful of thresholds, all selected by grid search on the validation split (never on test) to maximise fault macro-F1 plus load macro-F1.
\begin{enumerate}
    \item \textbf{Noise estimation:} $\hat{\sigma}_v = 1.4826\,\mathrm{MAD}(\Delta\tilde{v})/\sqrt{2}$ over $t/T\in[0.11,2.45]$ (robust to the few edges present), and $\hat{\sigma}_{dv}$ as the MAD of the derivative channel.
    \item \textbf{Edge detection:} peaks of $|d\tilde{v}|$ above $k_\sigma\hat{\sigma}_{dv}$ in the fault window $t/T\in[0.11,1.00]$. A peak at sample $p$ is accepted if its step $\hat{\Gamma} = \mathrm{med}\,\tilde{v}_s[p{+}3{:}p{+}13] - \mathrm{med}\,\tilde{v}_s[p{-}13{:}p{-}3]$ satisfies $|\hat{\Gamma}| \ge \max(0.02, c\,\hat{\sigma}_v)$. The first accepted edge is the fault candidate; none gives \texttt{none}.
    \item \textbf{Hard faults:} decided only from the level after the edge relative to the level before it, $\rho = L_{\mathrm{after}}/L_{\mathrm{before}}$ with $L_{\mathrm{after}}$ the median over $[t_e{+}0.05,\min(t_e{+}0.25,1)]$: $\rho<a\Rightarrow$ \texttt{short}, $\rho>b\Rightarrow$ \texttt{open}. The end-of-trace level is deliberately not used, because it is set by the load.
    \item \textbf{Soft faults by shape:} a persistent step down/up gives \texttt{shunt\_R}/\texttt{series\_R}; an excursion that returns to the pre-edge level within the window of a short section is a transient: a positive bump, or a dip with comparable falling and rising edges, gives \texttt{z0\_section}, and a sharp dip with slow recovery gives \texttt{shunt\_C}.
    \item \textbf{Location:} $\hat{d}/l = t_e/T - 0.05$, with $t_e$ the half-amplitude crossing of the edge.
    \item \textbf{Load (DC de-embedding):} from the settled end level $L_{\mathrm{end}}$ (median over $t/T\in[2.30,2.45]$) the input impedance is $Z_{\mathrm{in}} = Z_0 L_{\mathrm{end}}/(2-L_{\mathrm{end}})$. The estimated soft fault is removed (series: $Z_{\mathrm{in}}-R_f$; shunt: $(Z_{\mathrm{in}}^{-1}-R_f^{-1})^{-1}$, with $R_f$ from $\hat{\Gamma}$), then the line resistance $r_{dc}l$ is subtracted to give $\hat{R}_L$. The over/undershoot after the load echo distinguishes \texttt{RL}/\texttt{RC}; $|\Gamma_{dc}|$ small gives \texttt{matched}; $\hat{R}_L$ thresholds give \texttt{open}/\texttt{short}; $\hat{\tau}/T$ comes from a log-linear fit of the tail.
\end{enumerate}

\subsection{LightGBM with 94 Physics Features}
We engineer 94 features from the preprocessed 3-channel input:
\begin{itemize}
    \item \textbf{Coarse waveform (64):} bin averages of $\tilde{v}$ over 64 equal time bins.
    \item \textbf{Level statistics (7):} final level, minimum, maximum and their times for $t/T>0.11$, and mean levels over $[0.11,1.0]$ and $[1.0,2.5]$.
    \item \textbf{Edges (16):} for the four largest $|d\tilde{v}|$ peaks: time, signed peak height, local step size and half-maximum width.
    \item \textbf{First discontinuity (2):} first-edge reflection estimate $\hat{\Gamma}_1$ and impedance ratio $(1+\hat{\Gamma}_1)/(1-\hat{\Gamma}_1)$.
    \item \textbf{Load region (4):} level over $[0.90,0.97]$, overshoot and undershoot relative to the final level, and $\hat{\tau}/T$ from a log-linear tail fit.
    \item \textbf{Noise (1):} robust noise estimate $\hat{\sigma}$.
\end{itemize}
Four LightGBM \cite{ke2017lightgbm} models are trained: fault class (multiclass), location (regression on faulted samples), load class (multiclass where the load is observable) and load values ($\log_{10}R_L$ and $\log_{10}\tau/T$, each on its own mask). Each uses up to 3000 trees (learning rate 0.05, 63 leaves) with early stopping after 100 rounds on the validation set. Training uses two independent noisy copies of the training set (SNR drawn uniformly in [20, 50] dB per sample).

\subsection{Deep Multi-Task 1D-ResNet Architecture}
To exploit continuous temporal hierarchies and avoid manual feature design, we design a deep Multi-Task 1D-ResNet \cite{he2016deep} operating directly on the $[3, 1024]$ input tensor:
\begin{itemize}
    \item \textbf{Stem Layer:} A 1D convolution with kernel size 7, stride 2, padding 3 mapping $3 \to 32$ channels, followed by 1D Batch Normalization and ReLU activation, reducing sequence length to 512.
    \item \textbf{Four Residual Stages:} Each stage comprises two residual blocks with skip connections:
    \begin{itemize}
        \item Stage 1: 32 channels, stride 1 (length 512).
        \item Stage 2: 64 channels, stride 2 (length 256).
        \item Stage 3: 128 channels, stride 2 (length 128).
        \item Stage 4: 128 channels, stride 2 (length 64).
    \end{itemize}
    \item \textbf{Latent Embedding:} the $128\times64$ feature map is flattened (8192 values; flattening, unlike global pooling, keeps the absolute time position needed for fault location), followed by 20\% dropout, a linear projection $8192 \to 256$ and ReLU.
    \item \textbf{Task-Specific Prediction Heads:}
    \begin{itemize}
        \item Head 1 (T1 Fault): Linear($256 \to 7$) producing class logits $\mathbf{z}_{\mathrm{fault}}$.
        \item Head 2 (T2 Location): Linear($256 \to 1$) followed by Sigmoid activation, constraining $\hat{d}/l \in [0, 1]$.
        \item Head 3 (T3 Load Class): Linear($256 \to 6$) producing class logits $\mathbf{z}_{\mathrm{load}}$.
        \item Head 4 (T3 Load Values): Linear($256 \to 2$) regressing log-scaled parameters $[\log_{10} \hat{R}_L, \log_{10}(\hat{\tau}/T)]$.
    \end{itemize}
\end{itemize}

\subsection{Training Protocol}
The ResNet is trained with AdamW (learning rate $2\times10^{-3}$, weight decay $10^{-4}$, batch size 256) for at most 40 epochs, with a one-epoch linear warm-up followed by cosine decay, early stopping on the validation loss (patience 8), and mixed precision on GPU. Fresh measurement noise (SNR uniform in [20, 50] dB per sample) is drawn for every batch, so the network never sees the same noisy trace twice; noise injection and preprocessing run on the GPU. The validation set uses one fixed noisy copy at 30 dB. Channels are standardised with statistics from the training set. Main results are reported as mean $\pm$ standard deviation over three seeds.

\subsection{Physics-Masked Multi-Task Loss Formulation}
Joint multi-task optimization is governed by our physics-masked loss function:
\begin{equation}
\mathcal{L}_{\mathrm{total}} = w_1 \mathcal{L}_{\mathrm{fault}} + w_2 \mathcal{L}_{\mathrm{loc}} + w_3 \mathcal{L}_{\mathrm{load\_cls}} + w_4 \mathcal{L}_{\mathrm{load\_reg}}
\label{eq:total_loss}
\end{equation}
where loss balancing weights are set to $w_1 = 1.0, w_2 = 5.0, w_3 = 1.0, w_4 = 2.0$.

The component losses strictly respect the physical masking conditions:
\begin{enumerate}
    \item \textbf{Fault Classification Loss (Unmasked):}
    \begin{equation}
    \mathcal{L}_{\mathrm{fault}} = -\log\left( \frac{\exp(z_{\mathrm{fault}, y_1})}{\sum_{c=1}^7 \exp(z_{\mathrm{fault}, c})} \right)
    \end{equation}
    \item \textbf{Fault Location Loss (Masked on Fault Presence):} Evaluated only when true fault class is not \texttt{none} ($m_{\mathrm{loc}} = \mathbb{I}(y_1 \ne \text{\texttt{none}})$):
    \begin{equation}
    \mathcal{L}_{\mathrm{loc}} = \frac{\sum_{i=1}^B m_{\mathrm{loc}, i} \cdot \mathrm{Smooth}_{L_1}(\hat{d}_i/l - d_i/l)}{\sum_{i=1}^B m_{\mathrm{loc}, i} + \epsilon}
    \end{equation}
    \item \textbf{Load Classification Loss (Masked on Transmissibility):} Evaluated only when the fault is not a hard short or open ($m_{\mathrm{load}} = \mathbb{I}(y_1 \notin \{\text{\texttt{short}}, \text{\texttt{open}}\})$):
    \begin{equation}
    \mathcal{L}_{\mathrm{load\_cls}} = \frac{\sum_{i=1}^B m_{\mathrm{load}, i} \cdot \mathrm{CE}(\mathbf{z}_{\mathrm{load}, i}, y_{3, i})}{\sum_{i=1}^B m_{\mathrm{load}, i} + \epsilon}
    \end{equation}
    \item \textbf{Load Value Regression Loss (Masked):}
    \begin{equation}
    \mathcal{L}_{\mathrm{load\_reg}} = \frac{\sum_{i=1}^B m_{\mathrm{load}, i} \cdot \mathrm{Smooth}_{L_1}(\hat{\mathbf{r}}_i - \mathbf{r}_i)}{\sum_{i=1}^B m_{\mathrm{load}, i} + \epsilon}
    \end{equation}
\end{enumerate}
Masking guarantees that non-existent targets never inject corrupting gradients into the shared residual backbone.

